\gddchapter{Research objectives \& Participant information}


\section{Research objectives}

\subsection{Problem statement}

When people are  playing a really well made game some of them  describe feeling like they're directly inside of the digital experience.
Video games vary in their degree in which they can achieve that. This thesis explores weather the voice input modality can have a positive impact on said effect. 
While commands do exist in certain games but they're often require memorisation of the right commands.
These tests attempt to verify if deep integration with a large language model can possibly allow for sucessful 
navigation in a video game environment. 

\subsection{Research questions}

\begin{enumerate}
    \item How would the participants describe the experience of using voice input in the presented
game environment?
    \item How do users feel about using the voice modality compared to traditional input?
\end{enumerate}

\section{Ethical checklist}
    \begin{todolist}
    %   \item[\done] Informed consent
    \item Informed consent - The participant is consenting to the interview
    \item Recording consent - The user agrees to the interview being recorded
    \item Anonymity choice - real name or a pseudonym (reflected on the main page)
    \item Data storage - Confirm the data storage policy with the participant
    \end{todolist}

\section{Participant information}
\begin{table}[h]
\centering
\begin{tabular}{|p{5cm}|p{10cm}|}
\hline
\textbf{Field} & \textbf{Response} \\ \hline
Participant ID / Name & \\[1.5cm] \hline
Gaming Experience (e.g., Frequent, Casual, Non-gamer) & \\[1.5cm] \hline
Familiarity with Voice Assistants (e.g., Alexa, Siri) & \\[1.5cm] \hline
Physical Accessibility Context & \\[2cm] \hline
\end{tabular}
\end{table}
